\chapter{Operar con decimales}\label{ch:g5:decimalops}

Los números decimales suman, restan y multiplican casi como números enteros.
--- todo el arte es saber a dónde va el punto. este capítulo
cubre los métodos de columnas y la magia de multiplicar o dividir por
$10$, $100$, $1000$. (Multiplicar dos decimales juntos espera
\cref{ch:g6:decimals}.)

\section{Sumar y restar}

\begin{method}[Columnas con un punto]\label{met:g5:decimalops:addsub}
\begin{enumerate}
  \item Escribe los números con el \emph{puntos decimales alineado}
        (entonces las unidades están debajo de las unidades, las décimas debajo de las décimas);
  \item rellenar la parte decimal más corta con \omterm{def:g1:counting:zero}{ceros};
  \item sumar o restar como con números enteros;
  \item suelte el punto del resultado hacia abajo.
\end{enumerate}
\end{method}

\begin{example}\label{ex:g5:decimalops:addsub}
$27.5 + 3.86$ (almohadilla $27.5 = 27.50$):
\[
\begin{array}{r}
2\,7.5\,0 \\
+\ \ \ 3.8\,6 \\
\hline
3\,1.3\,6
\end{array}
\qquad\qquad
\begin{array}{r}
1\,2.4\,0 \\
-\ \ \ 5.6\,5 \\
\hline
\ \ \,6.7\,5
\end{array}
\]
Para la resta ($12.4 - 5.65$): rellenar, pedir prestado como de costumbre y comprobar
añadiendo de nuevo: $6.75 + 5.65 = 12.4$.
\end{example}

\begin{example}[Complementos mentales]\label{ex:g5:decimalops:complement}
¿Qué se debe agregar a $7.6$ alcanzar $10$? Trepar: $7.6 \to 8$ es
$0.4$, entonces $8 \to 10$ es $2$: el complemento es $2.4$. Útil para
cambiar: pagar $10$ por un $7.60$ artículo, obtener $2.40$ atrás.
\end{example}

\section{Diez veces más grande, diez veces más pequeño}

\begin{proposition}[Veces y dividido por 10, 100, 1000]\label{prop:g5:decimalops:shift}
Multiplicando por $10$ hace que cada dígito valga diez veces más: el
Los dígitos se deslizan un lugar hacia la izquierda --- que se parece al \emph{punto
moviendo un lugar a la derecha}. Dividir hace lo contrario:
\[
3.47 \times 10 = 34.7, \qquad
3.47 \times 100 = 347, \qquad
52.1 \div 10 = 5.21, \qquad
6 \div 100 = 0.06 .
\]
\end{proposition}
\admitted

\begin{omfigure}
\begin{tikzpicture}[scale=0.9]
  \draw (0,0) grid[xstep=1.5, ystep=0.75] (6.0,2.25);
  \node[font=\small] at (0.75,2.6) {tens};
  \node[font=\small] at (2.25,2.6) {units};
  \node[font=\small] at (3.75,2.6) {tenths};
  \node[font=\small] at (5.25,2.6) {hundredths};
  \node[font=\normalsize, omExo] at (2.25,1.85) {$3$};
  \node[font=\normalsize, omExo] at (3.75,1.85) {$4$};
  \node[font=\normalsize, omExo] at (5.25,1.85) {$7$};
  \node[font=\normalsize, omDef] at (0.75,1.1) {$3$};
  \node[font=\normalsize, omDef] at (2.25,1.1) {$4$};
  \node[font=\normalsize, omDef] at (3.75,1.1) {$7$};
  \node[right, font=\small] at (6.2,1.85) {$3.47$};
  \node[right, font=\small, omDef] at (6.2,1.1)
    {$3.47 \times 10 = 34.7$};
  \draw[-Stealth, omDef, thick] (4.4,1.7) -- (3.1,1.25);
\end{tikzpicture}

{\small Multiplicando por $10$: cada dígito sube un lugar hacia el
izquierda. El punto realmente no se mueve, los dígitos sí.}
\end{omfigure}

\begin{example}[Conversiones de nuevo]\label{ex:g5:decimalops:convert}
Las unidades de medida están separadas por potencias de diez, por lo que las conversiones son simplemente
estos turnos: $2.35$ metro $= 235$ centímetros ($\times 100$); $470$ gramo
$= 0.47$ kilos ($\div 1000$); $1.5$ l $= 150$ CL.
\end{example}

\section{Multiplicar un decimal por un número entero}

\begin{method}[Tiempos decimales enteros]\label{met:g5:decimalops:mult}
Calcule como si no tuviera sentido, luego proporcione el resultado tantos
dígitos decimales como factor decimal:
\[
4.35 \times 6: \quad 435 \times 6 = 2\,610
\ \to\ 4.35 \times 6 = 26.10 = 26.1 .
\]
Estima para ubicar el punto con confianza: $4.35$ esta cerca de
$4$, y $4 \times 6 = 24$ --- entonces $26.1$, no $2.61$ ni $261$.
\end{method}

\begin{example}[La suma repetida todavía funciona]\label{ex:g5:decimalops:mult}
$0.75 \times 4 = 0.75 + 0.75 + 0.75 + 0.75 = 3$: cuatro cuarteles hacer
un todo, en ropa decimal. Multiplicar por un número entero todavía es
``tantas veces''.
\end{example}

\begin{example}[Compras]\label{ex:g5:decimalops:shopping}
Tres cuadernos en $2.45$ cada uno y un bolígrafo en $1.20$:
\begin{enumerate}
  \item cuadernos: $2.45 \times 3 = 7.35$;
  \item total: $7.35 + 1.20 = 8.55$;
  \item cambiar de $10$: $10 - 8.55 = 1.45$.
\end{enumerate}
\end{example}

\section{Ejercicios}

\begin{exercise}[$\star $]\label{exo:g5:decimalops:1}
Calcular en columnas: $45.7 + 8.93$; \quad $16.25 + 3.75$; \quad
$0.86 + 12.4$.
\end{exercise}

\begin{exercise}[$\star $]\label{exo:g5:decimalops:2}
Calcule en columnas y verifique: $9.4 - 2.72$; \quad $20 - 13.45$;
\quad $6.03 - 5.9$.
\end{exercise}

\begin{exercise}[$\star $]\label{exo:g5:decimalops:3}
Complementos mentales a $10$ (como en
\cref{ex:g5:decimalops:complement}): $6.2$; \quad $3.75$; \quad
$9.99$.
\end{exercise}

\begin{exercise}[$\star $]\label{exo:g5:decimalops:4}
Calcular sin columnas:
\[
7.24 \times 10, \qquad
0.58 \times 100, \qquad
34.5 \div 10, \qquad
7 \div 100, \qquad
0.3 \times 1000 .
\]
\end{exercise}

\begin{exercise}[$\star $]\label{exo:g5:decimalops:5}
Convertir: $4.05$ m en centímetros; \quad $325$ centímetros en metros; \quad $0.6$ kg en g;
\quad $85$ cL en L.
\end{exercise}

\begin{exercise}[$\star $]\label{exo:g5:decimalops:6}
Primero estime y luego calcule: $6.2 \times 4$; \quad
$3.45 \times 8$; \quad $12.5 \times 6$.
\end{exercise}

\begin{exercise}[$\star $]\label{exo:g5:decimalops:7}
Una vuelta a una pista mide $0.4$ km. cuanto duran $7$ vueltas? Y
$25$ vueltas?
\end{exercise}

\begin{exercise}[$\star $]\label{exo:g5:decimalops:8}
lea compra $2$ baguettes en $1.15$ cada uno y un pastel en $12.60$. ella
paga con un $20$ factura. Calcule su total y su cambio.
\end{exercise}

\begin{exercise}[$\star $]\label{exo:g5:decimalops:9}
una botella sostiene $1.5$ L. ¿Cuántos litros hay en un paquete de $6$ botellas?
una familia bebe $0.75$ L por día: ¿por cuántos días dura un paquete?
último? (¿Cuántas veces $0.75$ encajar $9$?)
\end{exercise}

\begin{exercise}[$\star\star $]\label{exo:g5:decimalops:10}
Milo calcula $5.6 \times 3 = 15.18$, razonamiento ``$5 \times 3 = 15$
y $6 \times 3 = 18$''. Utilice una estimación para demostrar que la respuesta debe ser
incorrecto, luego calcule correctamente.
\end{exercise}

\begin{exercise}[$\star\star $]\label{exo:g5:decimalops:11}
Las cuatro vueltas de un corredor duraron $65.4$ s, $64.9$ s, $65.0$ arena
$66.1$ s.
\begin{enumerate}
  \item Calcula el tiempo total de las cuatro vueltas.
  \item El récord de cuatro vueltas es $260$ s. ¿Lo superó? por como
        mucho?
\end{enumerate}
\end{exercise}
